\begin{claim}{C.22}{Proved under stated hypotheses}{A countable scalar form suffices for the general raw adjoint}\label{C22}
Let $s\subseteq k$ be a countable pure unital subring, let $\mathcal A_s$ be a small $s$-linear category with projective $s$-module Hom groups, and put $\mathcal A_k=k\otimes_s\mathcal A_s$. Write
\[
 \mathcal M_r=\operatorname{Fun}_{r}(\mathcal A_r^{\mathrm{op}},\Mod r)
 \quad(r=s,k),
 \qquad L=k\otimes_s-:\mathcal M_s\rightleftarrows\mathcal M_k:U.
\]
Assume these two module categories have cocontinuous symmetric monoidal tensor products that are $r$-bilinear over their respective scalar rings $r=s,k$, with the following compatible scalar-form data:
\begin{enumerate}
\item the unit $\mathbf1_s$ has finitely generated projective values $\mathbf1_s(a)$;
\item there is a specified unit identification $L\mathbf1_s\cong\mathbf1_k$;
\item there are natural isomorphisms
\[
 U(Y\otimes_k LX)\cong UY\otimes_sX
 \quad(X\in\mathcal M_s,\ Y\in\mathcal M_k),
\]
compatible with units, associativity, and the scalar-adjunction unit $X\to ULX$.
\end{enumerate}
Here the subscripts on the internal tensor distinguish the two monoidal categories; they do not assert that their products are pointwise.

If $\mathcal M_k$ satisfies the manuscript's canonical normalization and full interval condition, then every projective resolving adjunction for $\varphi_k$ has a raw right adjoint inducing the desired equivalence. In particular a counterexample within the manuscript's corollary cannot admit such a countable scalar form. We do not assert that every $k$-linear category and tensor in the corollary has one.
\end{claim}
\begin{proof}
Let $T_s(a)=\operatorname{Hom}_s(\mathbf1_s(a),s)$. Objectwise finite projectivity gives
\[
 \mathbf1_s(a)\otimes_sV\cong\operatorname{Hom}_s(T_s(a),V).
\]
It follows by tensor--Hom adjunction for modules over a small linear category that the scalar-copower functor $\kappa_s(V)=V\otimes_s\mathbf1_s$ (the external, objectwise scalar copower) has left adjoint
\[
 \varphi_s(X)=T_s\otimes_{\mathcal A_s}X.
\]
The scalar-copower functor $\kappa_s$ is exact and preserves all limits, since its values are tensor products with finitely generated projective modules. Its usual right adjoint is $\operatorname{Hom}_{\mathcal M_s}(\mathbf1_s,-)$. The same statements hold over $k$, and finite-projective dual base change identifies $T_k=k\otimes_sT_s$. Thus (a), (b), and the additional module hypotheses needed in Claim~\ref{C6} hold over $s$; normalization and the interval will be checked below.

We choose the relative bar resolution of $T_s$. In degree $n$, it is a direct sum of covariant representables with coefficient modules
\[
 T_s(a_0)\otimes_s\mathcal A_s(a_0,a_1)
 \otimes_s\cdots\otimes_s\mathcal A_s(a_{n-1},a_n).
\]
All these coefficients are projective over $s$. The augmentation is split exact as a complex of $s$-modules at each object, using the insertion of its identity morphism in the usual bar contraction. Consequently this is a projective resolution in the covariant functor category. Its scalar extension is the corresponding bar resolution of $T_k$, even though $k$ is not assumed flat over $s$.

Let $\mathsf B^s$ and $\mathsf B^k$ be the resulting exact cocontinuous resolving functors. Their coend formulas give
\[
 \mathsf B^s(UY)\cong U\mathsf B^k(Y),
 \qquad
 \mathsf B^k(LX)\cong k\otimes_s\mathsf B^s(X).
 \tag{C-form-bar}\label{Cformbar}
\]
For the first identity, one moves the scalar-extension coefficient modules across the tensor products with the $k$-module values of $Y$; all remaining coend relations are exactly those for $\mathcal A_s$. The second is the corresponding base-change identity on both variables. These formulas do not use preservation of arbitrary kernels by scalar extension.

Given the interval $\alpha:\mathbf1_k\to J$ over $k$, put
\[
 \beta:\mathbf1_s\longrightarrow UL\mathbf1_s
 \cong U\mathbf1_k\longrightarrow UJ.
\]
For each $X\in\mathcal M_s$, the unit-compatible projection formula identifies $\beta\otimes_sX$ with
\[
 X\longrightarrow ULX
 \xrightarrow{\ U(\alpha\otimes_k LX)\ }
 U(J\otimes_k LX)\cong UJ\otimes_sX.
\]
The first map is pointwise monic because $s\to k$ is pure; the second is monic by (d). Its target is $L\varphi_s$-null: equation~\eqref{Cformbar} identifies its resolving complex with the restriction of the null complex $\mathsf B^k(J\otimes_kLX)$. This proves all of (d) over $s$.

For (c), apply \eqref{Cformbar} and the unit identification to the canonical augmentation $\mathsf B^s(\mathbf1_s)\to s[0]$. Its scalar extension is the canonical normalized augmentation over $k$. The pure-base-change homology injection proved in \eqref{Cpurehomology}, applied to its cone, implies that the augmentation over $s$ is a quasi-isomorphism. We have checked the actual canonical map, not only its graded homology groups.

We next arrange that the resolving functor sees every evaluation in degree zero. Let
\[
 E_s(X)=\bigoplus_{a\in\mathcal A_s}X(a),
 \qquad
 F^{(s)}=\mathsf B^s\oplus D^1E_s.
\]
The second summand is a contractible projective cocontinuous functor complex. Thus $F^{(s)}$ is still an exact projective resolution of $\varphi_s$, with an exact raw right adjoint $G^{(s)}$ by Claim~\ref{C6}. The full normalized hypotheses over $s$ now hold. The tensor is cocontinuous, hence the interval images are pure by Claim~\ref{C14}. Since $s$ is countable, that claim proves that this augmented raw adjunction is an equivalence after localization.

Base change the entire augmented resolution to obtain $F^{(k)}=\mathsf B^k\oplus D^1E_k$. Both \eqref{Cformbar} and the adjunction identity
\[
 U\operatorname{Hom}_{\Ch(k)}
      (k\otimes_s F^{(s)}h_a,K)
 \cong
 \operatorname{Hom}_{\Ch(s)}(F^{(s)}h_a,UK)
\]
show that the right adjoints and their counits agree after scalar restriction. Here $h_a$ denotes the appropriate representable over its indicated scalar ring. Therefore the raw adjunction $F^{(k)},G^{(k)}$ also has a quasi-isomorphism counit on every complex. As in Claim~\ref{C21}, this step requires no assumption that restriction of scalars preserves projective $k$-modules.

Let $R_a$ denote the right adjoint of evaluation on $\mathcal M_k$. Since $E_k$ is a direct summand of $F^{(k)}_0$, the product $\prod_aR_a(V)$ is a retract of $G^{(k)}(D^1V)$ for every module $V$. The latter is $L\varphi_k$-null because this augmented adjunction is already known to be homotopical. Hence $\prod_aR_a(V)$ is null for every $V$.

Finally take any originally chosen projective resolving adjunction $F\dashv G$. Each $F_n$ is a retract of a sum of evaluations, so $G$ on a product of disks is a retract of a product $\prod_tR_{a_t}(V_t)$. Exactly as in Claim~\ref{C21}, this is a retract of
\[
 \prod_a R_a\!\left(\prod_tV_t\right).
\]
The latter has just been proved null. The exact gate in Claim~\ref{C10} follows, and establishes the raw equivalence for the original $F,G$.

The countable scalar form is a substantive hypothesis. In particular, restricting scalars on the original $k$-linear category without descending its morphisms is insufficient: the resulting scalar-unit functor is generally $\kappa_k(k\otimes_s-)$, which need not preserve products. Nor is a countable subring assumed to contain an arbitrarily large set of composition or tensor structure constants. These are precisely the scope restrictions in the statement.
\end{proof}
